\documentclass[12pt,a4paper]{article}
\usepackage[utf8]{inputenc}
\usepackage[T1]{fontenc}
\usepackage[spanish]{babel}
\usepackage{amsmath,amssymb,amsfonts}
\usepackage{geometry}
\usepackage{booktabs}
\usepackage{hyperref}
\geometry{margin=2.5cm}

\title{Ontología de las Regiones Causalmente Auto-Consistentes}
\author{Kevin Mu\~noz}
\date{}

\begin{document}
\maketitle

\begin{abstract}
Desarrollamos una articulación ontológica compatible con un enfoque centrado en el horizonte para la cosmología y la física gravitacional. En lugar de introducir una nueva teoría física o proponer leyes dinámicas adicionales, el presente trabajo explora los compromisos ontológicos mínimos requeridos por marcos en los que la accesibilidad causal y la estructura del horizonte desempeñan un papel estructuralmente primario.

Dentro de esta perspectiva, introducimos la noción de región causalmente auto-consistente (RCAC): un dominio físico caracterizado por la accesibilidad causal y acotado por una estructura de horizonte que restringe los grados de libertad relevantes disponibles para observadores internos.

El concepto está destinado como una clarificación ontológica, no como un modelo completo de la realidad. No asume un sustrato microscópico específico, una completación global del espacio-tiempo ni una realización cosmológica única. En cambio, proporciona una estructura conceptual mínima capaz de organizar descripciones centradas en el horizonte de manera consistente con la termodinámica gravitacional y los marcos causales.

El interior de las geometrías de agujeros negros se discute únicamente como una realización ilustrativa y no como una identificación literal.
\end{abstract}

\section{Introducción}

Los desarrollos en termodinámica gravitacional y física de horizontes sugieren cada vez más que las fronteras causales poseen significancia estructural más allá de su interpretación geométrica tradicional. La termodinámica de agujeros negros, los horizontes cosmológicos y los enfoques basados en estructura causal indican consistentemente que la entropía, la accesibilidad a la información y las descripciones dependientes del observador pueden estar restringidas por horizontes, en lugar de estar determinadas únicamente por propiedades volumétricas.

Motivados por estas observaciones, se han propuesto marcos centrados en el horizonte en los que la estructura causal se trata como un principio organizador para la interpretación cosmológica.

El presente trabajo no extiende ni modifica tales marcos. En cambio, examina qué compromisos ontológicos mínimos se requieren si la accesibilidad causal y la estructura del horizonte se consideran físicamente significativas.

El propósito no es proporcionar una ontología fundamental de la realidad, sino articular un concepto de dominio capaz de apoyar descripciones centradas en el horizonte, permaneciendo compatible con múltiples realizaciones efectivas.

\section{Estatus y Alcance de las Afirmaciones}

El marco desarrollado aquí es de carácter ontológico y conceptual.

No introduce:

\begin{itemize}
  \item nuevas leyes dinámicas
  \item nuevos modelos cosmológicos
  \item descripciones microscópicas
  \item modificaciones de la Relatividad General
  \item predicciones observacionales
\end{itemize}

La noción de región causalmente auto-consistente debe, por tanto, interpretarse como una estructura de consistencia mínima, no como una teoría física independiente.

No se realiza ningún supuesto respecto a:

\begin{itemize}
  \item la totalidad del espacio-tiempo
  \item la existencia de dominios externos
  \item la discretización de los grados de libertad subyacentes
  \item mecanismos de gravedad cuántica
  \item reducción fundamental de la geometría
\end{itemize}

Las afirmaciones relativas a emergencia y ontología deben, por tanto, interpretarse al nivel de descripción estructural, no como afirmaciones metafísicas.

\section{Regiones Causalmente Auto-Consistentes}

\subsection{Definición Central}

Una región causalmente auto-consistente (RCAC) se define como el dominio físico mínimo requerido para una descripción consistente en términos de accesibilidad causal y estructura del horizonte.

Una RCAC se caracteriza por:

\begin{itemize}
  \item accesibilidad causal entre los grados de libertad internos
  \item una frontera definida por el horizonte
  \item observables operacionalmente significativos
  \item consistencia interna entre la información accesible y la dinámica efectiva
\end{itemize}

El concepto no implica que tales regiones constituyan la totalidad de la realidad física.

Más bien, representan los dominios mínimos dentro de los cuales las descripciones centradas en el horizonte se vuelven operacionalmente significativas.

No se realiza ningún supuesto respecto a la inclusión dentro de estructuras mayores.

En la realización efectiva cosmológica, el borde del RCAC se ubica en el radio del horizonte aparente $R_A = 1/H$, y las perturbaciones del borde corresponden a $\psi = \delta R_A / R_A$.

\subsection{Principios Ontológicos}

\paragraph{Primacía Causal:} Las relaciones causales se tratan como estructuralmente previas a la descripción geométrica. La geometría espacial se considera una representación efectiva de la accesibilidad causal, en lugar de una entidad ontológicamente fundamental. No se asume ningún mecanismo particular para tal emergencia.

\paragraph{Cierre Causal Operacional:} La influencia física observable no cruza el horizonte dentro del dominio de descripción. Las afirmaciones relativas a regiones externas se consideran indefinidas en lugar de falsas. Este principio debe interpretarse operacionalmente y no implica aislamiento metafísico absoluto.

\paragraph{Dominio Definido por el Horizonte:} El horizonte determina el dominio físicamente relevante disponible para observadores internos. El horizonte se trata, por tanto, como una estructura definitoria de la descripción, no meramente como una superficie geométrica. No se realiza ninguna afirmación respecto a la constitución microscópica de los grados de libertad del horizonte.

\paragraph{Entropía Asociada al Borde:} Se asume que la entropía relevante para el dominio está asociada principalmente a la estructura del horizonte. Este principio sigue los conocimientos establecidos de la termodinámica gravitacional y debe entenderse como una guía estructural, no como una derivación. No se introduce ningún modelo cuantitativo de entropía.

\paragraph{Restricción de Consistencia:} Cualquier descripción efectiva admisible de una RCAC debe permanecer compatible con las propiedades causales y entrópicas de su frontera definitoria. Las descripciones de bulk se interpretan, por tanto, como restringidas por la estructura del borde. No se implica ninguna realización dinámica única.

En la realización cosmológica plana, esta restricción se ejemplifica con la relación $\delta\rho/\rho = -2\psi$: las perturbaciones de densidad en el bulk están determinadas por las fluctuaciones fraccionales de la frontera del horizonte, haciendo explícito cómo la estructura del borde restringe las descripciones del bulk.

\section{Emergencia del Espacio y el Tiempo}

Dentro de los marcos centrados en el horizonte, la geometría y el ordenamiento temporal no necesitan considerarse fundamentales.

Las relaciones espaciales pueden interpretarse, en cambio, como descripciones efectivas de la accesibilidad causal entre grados de libertad.

Del mismo modo, el ordenamiento temporal puede surgir de cambios en la información disponible y de relaciones causales en evolución.

El presente trabajo no intenta derivar estas estructuras.

Más bien, identifica compatibilidad conceptual con enfoques en los que el espacio-tiempo es emergente.

\section{Realizaciones y Ejemplos Ilustrativos}

El concepto de RCAC no prescribe una realización única del espacio-tiempo.

Diferentes contextos físicos pueden instanciar estructuras compatibles, incluyendo:

\begin{itemize}
  \item sistemas de horizonte cosmológico
  \item fases dominadas por el horizonte
  \item interiores de agujeros negros
  \item construcciones basadas en estructura causal
\end{itemize}

Estos ejemplos deben entenderse como realizaciones ilustrativas, no como casos definitorios.

La compatibilidad no implica equivalencia.

\section{Caso Ilustrativo: Geometría Interior de Agujeros Negros}

La región interior de un espacio-tiempo de agujero negro proporciona un ejemplo conceptual particularmente útil.

Tales sistemas poseen:

\begin{itemize}
  \item fronteras causales
  \item descripciones dependientes del observador
  \item entropía asociada al horizonte
  \item accesibilidad restringida
\end{itemize}

Estas características se asemejan a las requeridas de una región causalmente auto-consistente.

Esta comparación está destinada únicamente como analogía conceptual.

No se propone ninguna identificación entre dominios cosmológicos e interiores de agujeros negros.

Ninguna conclusión respecto al origen cosmológico se desprende de este ejemplo.

\section{Direcciones Abiertas}

Varias preguntas permanecen abiertas:

\begin{itemize}
  \item origen microscópico de los grados de libertad del horizonte
  \item relaciones cuantitativas borde-a-bulk
  \item emergencia de geometría efectiva
  \item relación con la gravedad cuántica
  \item posibles realizaciones efectivas
\end{itemize}

Estas preguntas abiertas definen direcciones de investigación futura, no deficiencias.

\section{Conclusión}

Hemos introducido la noción de regiones causalmente auto-consistentes como una articulación ontológica mínima compatible con los enfoques centrados en el horizonte.

El propósito del marco no es proporcionar una ontología completa ni reemplazar las descripciones físicas establecidas.

En cambio, ofrece una estructura conceptual mínima capaz de apoyar interpretaciones organizadas en torno a la accesibilidad causal, la estructura del horizonte y la entropía asociada al borde.

Los desarrollos futuros determinarán si estas ideas admiten realización cuantitativa dentro de marcos teóricos más amplios.

\end{document}
